
\documentclass{article}
\usepackage{colortbl}
\usepackage{makecell}
\usepackage{multirow}
\usepackage{supertabular}

\begin{document}

\newcounter{utterance}

\centering \large Interaction Transcript for game `dond', experiment `coop\_en', episode 1 with Qwen3{-}32B{-}t0.0.
\vspace{24pt}

{ \footnotesize  \setcounter{utterance}{1}
\setlength{\tabcolsep}{0pt}
\begin{supertabular}{c@{$\;$}|p{.15\linewidth}@{}p{.15\linewidth}p{.15\linewidth}p{.15\linewidth}p{.15\linewidth}p{.15\linewidth}}
    \# & $\;$A & \multicolumn{4}{c}{Game Master} & $\;\:$B\\
    \hline

    \theutterance \stepcounter{utterance}  
    & & \multicolumn{4}{p{0.6\linewidth}}{
        \cellcolor[rgb]{0.9,0.9,0.9}{
            \makecell[{{p{\linewidth}}}]{
                \texttt{\tiny{[P1$\langle$GM]}}
                \texttt{You are playing a cooperative negotiation game in which you have to agree on how to divide a set of items among you and another player.} \\
\\ 
\texttt{Rules:} \\
\texttt{(a) You and the other player are given a set of items. Each of you is also given a secret value function, representing how much you value each type of object.} \\
\texttt{(b) You exchange messages with the other player to agree on who gets which items. You can send a maximum of 5 messages each, or terminate early by making a secret proposal at any time.} \\
\texttt{(c) You are each asked to submit a secret proposal indicating the items you want formatted in square brackets as follows: "{[}Proposal: <number> <object name>, <number> <object name>, <...>{]}"} \\
\texttt{(d) If your proposals are complementary, i.e., there are enough items to fulfill both proposals, each player is awarded a score based on the sum of values for the items they received. Otherwise, both of you get zero points.} \\
\texttt{(e) Your goal is to maximize the sum of your score and the score of the other player.} \\
\\ 
\texttt{Let us start.} \\
\\ 
\texttt{The set of available items is:} \\
\\ 
\texttt{2 bells, 1 doll, 2 threads} \\
\\ 
\texttt{Your secret value function is:} \\
\\ 
\texttt{bell: 1, doll: 2, thread: 3} \\
\\ 
\texttt{IMPORTANT: Your messages, unless it is the secret proposal, are directly transmitted to the other player, so do not include any response to the rules or text announcing your message. To make a secret proposal, use the indicated format. Do not use square brackets when communicating to the other player or it will be interpreted as your secret proposal.} \\
            }
        }
    }
    & \\ \\

    \theutterance \stepcounter{utterance}  
    & \multicolumn{4}{p{0.6\linewidth}}{
        \cellcolor[rgb]{0.9,0.9,0.9}{
            \makecell[{{p{\linewidth}}}]{
                \texttt{\tiny{[P1$\rangle$GM]}}
                \texttt{<think>} \\
\texttt{Okay, let's see. I need to divide the items between me and another player. The items are 2 bells, 1 doll, and 2 threads. My value function is bell:1, doll:2, thread:3. The other player's value function is secret, so I don't know what they value. The goal is to maximize the total score for both of us.} \\
\\ 
\texttt{First, I should figure out what the other player might value. Since I don't have their value function, I need to make some assumptions or ask them. But since I can only send 5 messages, I need to be efficient. Maybe I can ask them what they value most, but they might not tell me the truth. Alternatively, I can propose a division that's fair based on possible scenarios.} \\
\\ 
\texttt{Let me think about possible value functions for the other player. They could value bells, dolls, or threads differently. For example, if they value threads highly, they might want more threads. If they value the doll more, they might prefer that. Since I have threads as 3 each, which is higher than bells and doll, maybe I should consider giving them threads if they don't value them as much. But without knowing their values, it's tricky.} \\
\\ 
\texttt{The key is to make a proposal that allows both of us to get items we value, so the total is maximized. Let's consider possible allocations.} \\
\\ 
\texttt{If I take the doll (2 points) and one thread (3), that's 5 points. The other player could take the remaining bell (1) and thread (3), totaling 4. But maybe there's a better split. Alternatively, if I take two threads (6 points) and the other takes the doll (2) and a bell (1), total would be 6 + 3 = 9. Wait, but the other player's value for the doll and bell depends on their own function. If they value the doll higher, maybe that's better.} \\
\\ 
\texttt{But since I don't know their values, I need to make a proposal that's likely to be complementary. Maybe propose taking the doll and one thread, leaving the bells and another thread for them. Alternatively, take two threads and a bell, but that might not be optimal.} \\
\\ 
\texttt{Alternatively, maybe I should ask the other player what they prefer. For example, I could message them: "Would you prefer the doll or the threads? Let me know so we can} \\
            }
        }
    }
    & & \\ \\

    \theutterance \stepcounter{utterance}  
    & & & \multicolumn{2}{p{0.3\linewidth}}{
        \cellcolor[rgb]{0.9,0.9,0.9}{
            \makecell[{{p{\linewidth}}}]{
                \texttt{\tiny{[GM$|$GM]}}
                \texttt{continue} \\
            }
        }
    }
    & & \\ \\

    \theutterance \stepcounter{utterance}  
    & & & \multicolumn{2}{p{0.3\linewidth}}{
        \cellcolor[rgb]{0.9,0.9,0.9}{
            \makecell[{{p{\linewidth}}}]{
                \texttt{\tiny{[GM$|$GM]}}
                \texttt{<think>} \\
\texttt{Okay, let's see. I need to divide the items between me and another player. The items are 2 bells, 1 doll, and 2 threads. My value function is bell:1, doll:2, thread:3. The other player's value function is secret, so I don't know what they value. The goal is to maximize the total score for both of us.} \\
\\ 
\texttt{First, I should figure out what the other player might value. Since I don't have their value function, I need to make some assumptions or ask them. But since I can only send 5 messages, I need to be efficient. Maybe I can ask them what they value most, but they might not tell me the truth. Alternatively, I can propose a division that's fair based on possible scenarios.} \\
\\ 
\texttt{Let me think about possible value functions for the other player. They could value bells, dolls, or threads differently. For example, if they value threads highly, they might want more threads. If they value the doll more, they might prefer that. Since I have threads as 3 each, which is higher than bells and doll, maybe I should consider giving them threads if they don't value them as much. But without knowing their values, it's tricky.} \\
\\ 
\texttt{The key is to make a proposal that allows both of us to get items we value, so the total is maximized. Let's consider possible allocations.} \\
\\ 
\texttt{If I take the doll (2 points) and one thread (3), that's 5 points. The other player could take the remaining bell (1) and thread (3), totaling 4. But maybe there's a better split. Alternatively, if I take two threads (6 points) and the other takes the doll (2) and a bell (1), total would be 6 + 3 = 9. Wait, but the other player's value for the doll and bell depends on their own function. If they value the doll higher, maybe that's better.} \\
\\ 
\texttt{But since I don't know their values, I need to make a proposal that's likely to be complementary. Maybe propose taking the doll and one thread, leaving the bells and another thread for them. Alternatively, take two threads and a bell, but that might not be optimal.} \\
\\ 
\texttt{Alternatively, maybe I should ask the other player what they prefer. For example, I could message them: "Would you prefer the doll or the threads? Let me know so we can} \\
            }
        }
    }
    & & \\ \\

    \theutterance \stepcounter{utterance}  
    & & \multicolumn{4}{p{0.6\linewidth}}{
        \cellcolor[rgb]{0.9,0.9,0.9}{
            \makecell[{{p{\linewidth}}}]{
                \texttt{\tiny{[GM$\rangle$P2]}}
                \texttt{You are playing a cooperative negotiation game in which you have to agree on how to divide a set of items among you and another player.} \\
\\ 
\texttt{Rules:} \\
\texttt{(a) You and the other player are given a set of items. Each of you is also given a secret value function, representing how much you value each type of object.} \\
\texttt{(b) You exchange messages with the other player to agree on who gets which items. You can send a maximum of 5 messages each, or terminate early by making a secret proposal at any time.} \\
\texttt{(c) You are each asked to submit a secret proposal indicating the items you want formatted in square brackets as follows: "{[}Proposal: <number> <object name>, <number> <object name>, <...>{]}"} \\
\texttt{(d) If your proposals are complementary, i.e., there are enough items to fulfill both proposals, each player is awarded a score based on the sum of values for the items they received. Otherwise, both of you get zero points.} \\
\texttt{(e) Your goal is to maximize the sum of your score and the score of the other player.} \\
\\ 
\texttt{Let us start.} \\
\\ 
\texttt{The set of available items is:} \\
\\ 
\texttt{2 bells, 1 doll, 2 threads} \\
\\ 
\texttt{Your secret value function is:} \\
\\ 
\texttt{bell: 0, doll: 2, thread: 4} \\
\\ 
\texttt{IMPORTANT: Your messages, unless it is the secret proposal, are directly transmitted to the other player, so do not include any response to the rules or text announcing your message. To make a secret proposal, use the indicated format. Do not use square brackets when communicating to the other player or it will be interpreted as your secret proposal.} \\
\\ 
\\ 
\texttt{<think>} \\
\texttt{Okay, let's see. I need to divide the items between me and another player. The items are 2 bells, 1 doll, and 2 threads. My value function is bell:1, doll:2, thread:3. The other player's value function is secret, so I don't know what they value. The goal is to maximize the total score for both of us.} \\
\\ 
\texttt{First, I should figure out what the other player might value. Since I don't have their value function, I need to make some assumptions or ask them. But since I can only send 5 messages, I need to be efficient. Maybe I can ask them what they value most, but they might not tell me the truth. Alternatively, I can propose a division that's fair based on possible scenarios.} \\
\\ 
\texttt{Let me think about possible value functions for the other player. They could value bells, dolls, or threads differently. For example, if they value threads highly, they might want more threads. If they value the doll more, they might prefer that. Since I have threads as 3 each, which is higher than bells and doll, maybe I should consider giving them threads if they don't value them as much. But without knowing their values, it's tricky.} \\
\\ 
\texttt{The key is to make a proposal that allows both of us to get items we value, so the total is maximized. Let's consider possible allocations.} \\
\\ 
\texttt{If I take the doll (2 points) and one thread (3), that's 5 points. The other player could take the remaining bell (1) and thread (3), totaling 4. But maybe there's a better split. Alternatively, if I take two threads (6 points) and the other takes the doll (2) and a bell (1), total would be 6 + 3 = 9. Wait, but the other player's value for the doll and bell depends on their own function. If they value the doll higher, maybe that's better.} \\
\\ 
\texttt{But since I don't know their values, I need to make a proposal that's likely to be complementary. Maybe propose taking the doll and one thread, leaving the bells and another thread for them. Alternatively, take two threads and a bell, but that might not be optimal.} \\
\\ 
\texttt{Alternatively, maybe I should ask the other player what they prefer. For example, I could message them: "Would you prefer the doll or the threads? Let me know so we can} \\
            }
        }
    }
    & \\ \\

    \theutterance \stepcounter{utterance}  
    & & & \multicolumn{4}{p{0.6\linewidth}}{
        \cellcolor[rgb]{0.9,0.9,0.9}{
            \makecell[{{p{\linewidth}}}]{
                \texttt{\tiny{[GM$\langle$P2]}}
                \texttt{<think>} \\
\texttt{Okay, let's see. I need to figure out how to divide the items between me and the other player to maximize the total score. The items are 2 bells, 1 doll, and 2 threads. My value function is bell: 0, doll: 2, thread: 4. The other player's values are secret, so I don't know what they value. The goal is to maximize the sum of both our scores.} \\
\\ 
\texttt{First, I should consider what the other player might value. Since I don't have their value function, I need to make some assumptions or ask them. But since I can only send 5 messages, I need to be efficient. Maybe I can ask them what they value most, but they might not tell me the truth. Alternatively, I can propose a division that's fair based on possible scenarios.} \\
\\ 
\texttt{Let me think about possible value functions for the other player. They could value bells, dolls, or threads differently. For example, if they value threads highly, they might want more threads. If they value the doll more, they might prefer that. Since I have threads as 4 each, which is higher than bells (0) and doll (2), maybe I should consider giving them threads if they don't value them as much. But without knowing their values, it's tricky.} \\
\\ 
\texttt{The key is to make a proposal that allows both of us to get items we value, so the total is maximized. Let's consider possible allocations.} \\
\\ 
\texttt{If I take the doll (2 points) and one thread (4), that's 6 points. The other player could take the remaining bell (0) and thread (4), totaling 4. But maybe there's a better split. Alternatively, if I take two threads (8 points) and the other takes the doll (2) and a bell (0), total would be 8 + 2 = 10. Wait, but the other player's value for the doll and bell depends on their own function. If they value the doll higher, maybe that's better.} \\
\\ 
\texttt{But since I don't know their values, I need to make a proposal that's likely to be complementary. Maybe propose taking the doll and one thread, leaving the bells and another thread for them. Alternatively, take two threads and a bell, but that might not be optimal.} \\
\\ 
\texttt{Alternatively, maybe I should ask the other player what they prefer. For example, I could message them:} \\
            }
        }
    }
     \\ \\

    \theutterance \stepcounter{utterance}  
    & & & \multicolumn{2}{p{0.3\linewidth}}{
        \cellcolor[rgb]{0.9,0.9,0.9}{
            \makecell[{{p{\linewidth}}}]{
                \texttt{\tiny{[GM$|$GM]}}
                \texttt{continue} \\
            }
        }
    }
    & & \\ \\

    \theutterance \stepcounter{utterance}  
    & & & \multicolumn{2}{p{0.3\linewidth}}{
        \cellcolor[rgb]{0.9,0.9,0.9}{
            \makecell[{{p{\linewidth}}}]{
                \texttt{\tiny{[GM$|$GM]}}
                \texttt{<think>} \\
\texttt{Okay, let's see. I need to figure out how to divide the items between me and the other player to maximize the total score. The items are 2 bells, 1 doll, and 2 threads. My value function is bell: 0, doll: 2, thread: 4. The other player's values are secret, so I don't know what they value. The goal is to maximize the sum of both our scores.} \\
\\ 
\texttt{First, I should consider what the other player might value. Since I don't have their value function, I need to make some assumptions or ask them. But since I can only send 5 messages, I need to be efficient. Maybe I can ask them what they value most, but they might not tell me the truth. Alternatively, I can propose a division that's fair based on possible scenarios.} \\
\\ 
\texttt{Let me think about possible value functions for the other player. They could value bells, dolls, or threads differently. For example, if they value threads highly, they might want more threads. If they value the doll more, they might prefer that. Since I have threads as 4 each, which is higher than bells (0) and doll (2), maybe I should consider giving them threads if they don't value them as much. But without knowing their values, it's tricky.} \\
\\ 
\texttt{The key is to make a proposal that allows both of us to get items we value, so the total is maximized. Let's consider possible allocations.} \\
\\ 
\texttt{If I take the doll (2 points) and one thread (4), that's 6 points. The other player could take the remaining bell (0) and thread (4), totaling 4. But maybe there's a better split. Alternatively, if I take two threads (8 points) and the other takes the doll (2) and a bell (0), total would be 8 + 2 = 10. Wait, but the other player's value for the doll and bell depends on their own function. If they value the doll higher, maybe that's better.} \\
\\ 
\texttt{But since I don't know their values, I need to make a proposal that's likely to be complementary. Maybe propose taking the doll and one thread, leaving the bells and another thread for them. Alternatively, take two threads and a bell, but that might not be optimal.} \\
\\ 
\texttt{Alternatively, maybe I should ask the other player what they prefer. For example, I could message them:} \\
            }
        }
    }
    & & \\ \\

    \theutterance \stepcounter{utterance}  
    & & \multicolumn{4}{p{0.6\linewidth}}{
        \cellcolor[rgb]{0.9,0.9,0.9}{
            \makecell[{{p{\linewidth}}}]{
                \texttt{\tiny{[P1$\langle$GM]}}
                \texttt{<think>} \\
\texttt{Okay, let's see. I need to figure out how to divide the items between me and the other player to maximize the total score. The items are 2 bells, 1 doll, and 2 threads. My value function is bell: 0, doll: 2, thread: 4. The other player's values are secret, so I don't know what they value. The goal is to maximize the sum of both our scores.} \\
\\ 
\texttt{First, I should consider what the other player might value. Since I don't have their value function, I need to make some assumptions or ask them. But since I can only send 5 messages, I need to be efficient. Maybe I can ask them what they value most, but they might not tell me the truth. Alternatively, I can propose a division that's fair based on possible scenarios.} \\
\\ 
\texttt{Let me think about possible value functions for the other player. They could value bells, dolls, or threads differently. For example, if they value threads highly, they might want more threads. If they value the doll more, they might prefer that. Since I have threads as 4 each, which is higher than bells (0) and doll (2), maybe I should consider giving them threads if they don't value them as much. But without knowing their values, it's tricky.} \\
\\ 
\texttt{The key is to make a proposal that allows both of us to get items we value, so the total is maximized. Let's consider possible allocations.} \\
\\ 
\texttt{If I take the doll (2 points) and one thread (4), that's 6 points. The other player could take the remaining bell (0) and thread (4), totaling 4. But maybe there's a better split. Alternatively, if I take two threads (8 points) and the other takes the doll (2) and a bell (0), total would be 8 + 2 = 10. Wait, but the other player's value for the doll and bell depends on their own function. If they value the doll higher, maybe that's better.} \\
\\ 
\texttt{But since I don't know their values, I need to make a proposal that's likely to be complementary. Maybe propose taking the doll and one thread, leaving the bells and another thread for them. Alternatively, take two threads and a bell, but that might not be optimal.} \\
\\ 
\texttt{Alternatively, maybe I should ask the other player what they prefer. For example, I could message them:} \\
            }
        }
    }
    & \\ \\

    \theutterance \stepcounter{utterance}  
    & \multicolumn{4}{p{0.6\linewidth}}{
        \cellcolor[rgb]{0.9,0.9,0.9}{
            \makecell[{{p{\linewidth}}}]{
                \texttt{\tiny{[P1$\rangle$GM]}}
                \texttt{<think>} \\
\texttt{Okay, let's see. I need to figure out how to divide the items between me and the other player to maximize the total score. The items are 2 bells, 1 doll, and 2 threads. My value function is bell: 0, doll: 2, thread: 4. The other player's values are secret, so I don't know what they value. The goal is to maximize the sum of both our scores.} \\
\\ 
\texttt{First, I should consider what the other player might value. Since I don't have their value function, I need to make some assumptions or ask them. But since I can only send 5 messages, I need to be efficient. Maybe I can ask them what they value most, but they might not tell me the truth. Alternatively, I can propose a division that's fair based on possible scenarios.} \\
\\ 
\texttt{Let me think about possible value functions for the other player. They could value bells, dolls, or threads differently. For example, if they value threads highly, they might want more threads. If they value the doll more, they might prefer that. Since I have threads as 4 each, which is higher than bells (0) and doll (2), maybe I should consider giving them threads if they don't value them as much. But without knowing their values, it's tricky.} \\
\\ 
\texttt{The key is to make a proposal that allows both of us to get items we value, so the total is maximized. Let's consider possible allocations.} \\
\\ 
\texttt{If I take the doll (2 points) and one thread (4), that's 6 points. The other player could take the remaining bell (0) and thread (4), totaling 4. But maybe there's a better split. Alternatively, if I take two threads (8 points) and the other takes the doll (2) and a bell (0), total would be 8 + 2 = 10. Wait, but the other player's value for the doll and bell depends on their own function. If they value the doll higher, maybe that's better.} \\
\\ 
\texttt{But since I don't know their values, I need to make a proposal that's likely to be complementary. Maybe propose taking the doll and one thread, leaving the bells and another thread for them. Alternatively, take two threads and a bell, but that might not be optimal.} \\
\\ 
\texttt{Alternatively, maybe I should ask the other player what they prefer. For example, I could message them:} \\
            }
        }
    }
    & & \\ \\

    \theutterance \stepcounter{utterance}  
    & & & \multicolumn{2}{p{0.3\linewidth}}{
        \cellcolor[rgb]{0.9,0.9,0.9}{
            \makecell[{{p{\linewidth}}}]{
                \texttt{\tiny{[GM$|$GM]}}
                \texttt{continue} \\
            }
        }
    }
    & & \\ \\

    \theutterance \stepcounter{utterance}  
    & & & \multicolumn{2}{p{0.3\linewidth}}{
        \cellcolor[rgb]{0.9,0.9,0.9}{
            \makecell[{{p{\linewidth}}}]{
                \texttt{\tiny{[GM$|$GM]}}
                \texttt{<think>} \\
\texttt{Okay, let's see. I need to figure out how to divide the items between me and the other player to maximize the total score. The items are 2 bells, 1 doll, and 2 threads. My value function is bell: 0, doll: 2, thread: 4. The other player's values are secret, so I don't know what they value. The goal is to maximize the sum of both our scores.} \\
\\ 
\texttt{First, I should consider what the other player might value. Since I don't have their value function, I need to make some assumptions or ask them. But since I can only send 5 messages, I need to be efficient. Maybe I can ask them what they value most, but they might not tell me the truth. Alternatively, I can propose a division that's fair based on possible scenarios.} \\
\\ 
\texttt{Let me think about possible value functions for the other player. They could value bells, dolls, or threads differently. For example, if they value threads highly, they might want more threads. If they value the doll more, they might prefer that. Since I have threads as 4 each, which is higher than bells (0) and doll (2), maybe I should consider giving them threads if they don't value them as much. But without knowing their values, it's tricky.} \\
\\ 
\texttt{The key is to make a proposal that allows both of us to get items we value, so the total is maximized. Let's consider possible allocations.} \\
\\ 
\texttt{If I take the doll (2 points) and one thread (4), that's 6 points. The other player could take the remaining bell (0) and thread (4), totaling 4. But maybe there's a better split. Alternatively, if I take two threads (8 points) and the other takes the doll (2) and a bell (0), total would be 8 + 2 = 10. Wait, but the other player's value for the doll and bell depends on their own function. If they value the doll higher, maybe that's better.} \\
\\ 
\texttt{But since I don't know their values, I need to make a proposal that's likely to be complementary. Maybe propose taking the doll and one thread, leaving the bells and another thread for them. Alternatively, take two threads and a bell, but that might not be optimal.} \\
\\ 
\texttt{Alternatively, maybe I should ask the other player what they prefer. For example, I could message them:} \\
            }
        }
    }
    & & \\ \\

    \theutterance \stepcounter{utterance}  
    & & \multicolumn{4}{p{0.6\linewidth}}{
        \cellcolor[rgb]{0.9,0.9,0.9}{
            \makecell[{{p{\linewidth}}}]{
                \texttt{\tiny{[GM$\rangle$P2]}}
                \texttt{<think>} \\
\texttt{Okay, let's see. I need to figure out how to divide the items between me and the other player to maximize the total score. The items are 2 bells, 1 doll, and 2 threads. My value function is bell: 0, doll: 2, thread: 4. The other player's values are secret, so I don't know what they value. The goal is to maximize the sum of both our scores.} \\
\\ 
\texttt{First, I should consider what the other player might value. Since I don't have their value function, I need to make some assumptions or ask them. But since I can only send 5 messages, I need to be efficient. Maybe I can ask them what they value most, but they might not tell me the truth. Alternatively, I can propose a division that's fair based on possible scenarios.} \\
\\ 
\texttt{Let me think about possible value functions for the other player. They could value bells, dolls, or threads differently. For example, if they value threads highly, they might want more threads. If they value the doll more, they might prefer that. Since I have threads as 4 each, which is higher than bells (0) and doll (2), maybe I should consider giving them threads if they don't value them as much. But without knowing their values, it's tricky.} \\
\\ 
\texttt{The key is to make a proposal that allows both of us to get items we value, so the total is maximized. Let's consider possible allocations.} \\
\\ 
\texttt{If I take the doll (2 points) and one thread (4), that's 6 points. The other player could take the remaining bell (0) and thread (4), totaling 4. But maybe there's a better split. Alternatively, if I take two threads (8 points) and the other takes the doll (2) and a bell (0), total would be 8 + 2 = 10. Wait, but the other player's value for the doll and bell depends on their own function. If they value the doll higher, maybe that's better.} \\
\\ 
\texttt{But since I don't know their values, I need to make a proposal that's likely to be complementary. Maybe propose taking the doll and one thread, leaving the bells and another thread for them. Alternatively, take two threads and a bell, but that might not be optimal.} \\
\\ 
\texttt{Alternatively, maybe I should ask the other player what they prefer. For example, I could message them:} \\
            }
        }
    }
    & \\ \\

    \theutterance \stepcounter{utterance}  
    & & & \multicolumn{4}{p{0.6\linewidth}}{
        \cellcolor[rgb]{0.9,0.9,0.9}{
            \makecell[{{p{\linewidth}}}]{
                \texttt{\tiny{[GM$\langle$P2]}}
                \texttt{<think>} \\
\texttt{Okay, let's see. I need to figure out how to divide the items between me and the other player to maximize the total score. The items are 2 bells, 1 doll, and 2 threads. My value function is bell: 0, doll: 2, thread: 4. The other player's values are secret, so I don't know what they value. The goal is to maximize the sum of both our scores.} \\
\\ 
\texttt{First, I should consider what the other player might value. Since I don't have their value function, I need to make some assumptions or ask them. But since I can only send 5 messages, I need to be efficient. Maybe I can ask them what they value most, but they might not tell me the truth. Alternatively, I can propose a division that's fair based on possible scenarios.} \\
\\ 
\texttt{Let me think about possible value functions for the other player. They could value bells, dolls, or threads differently. For example, if they value threads highly, they might want more threads. If they value the doll more, they might prefer that. Since I have threads as 4 each, which is higher than bells (0) and doll (2), maybe I should consider giving them threads if they don't value them as much. But without knowing their values, it's tricky.} \\
\\ 
\texttt{The key is to make a proposal that allows both of us to get items we value, so the total is maximized. Let's consider possible allocations.} \\
\\ 
\texttt{If I take the doll (2 points) and one thread (4), that's 6 points. The other player could take the remaining bell (0) and thread (4), totaling 4. But maybe there's a better split. Alternatively, if I take two threads (8 points) and the other takes the doll (2) and a bell (0), total would be 8 + 2 = 10. Wait, but the other player's value for the doll and bell depends on their own function. If they value the doll higher, maybe that's better.} \\
\\ 
\texttt{But since I don't know their values, I need to make a proposal that's likely to be complementary. Maybe propose taking the doll and one thread, leaving the bells and another thread for them. Alternatively, take two threads and a bell, but that might not be optimal.} \\
\\ 
\texttt{Alternatively, maybe I should ask the other player what they prefer. For example, I could message them:} \\
            }
        }
    }
     \\ \\

    \theutterance \stepcounter{utterance}  
    & & & \multicolumn{2}{p{0.3\linewidth}}{
        \cellcolor[rgb]{0.9,0.9,0.9}{
            \makecell[{{p{\linewidth}}}]{
                \texttt{\tiny{[GM$|$GM]}}
                \texttt{continue} \\
            }
        }
    }
    & & \\ \\

    \theutterance \stepcounter{utterance}  
    & & & \multicolumn{2}{p{0.3\linewidth}}{
        \cellcolor[rgb]{0.9,0.9,0.9}{
            \makecell[{{p{\linewidth}}}]{
                \texttt{\tiny{[GM$|$GM]}}
                \texttt{<think>} \\
\texttt{Okay, let's see. I need to figure out how to divide the items between me and the other player to maximize the total score. The items are 2 bells, 1 doll, and 2 threads. My value function is bell: 0, doll: 2, thread: 4. The other player's values are secret, so I don't know what they value. The goal is to maximize the sum of both our scores.} \\
\\ 
\texttt{First, I should consider what the other player might value. Since I don't have their value function, I need to make some assumptions or ask them. But since I can only send 5 messages, I need to be efficient. Maybe I can ask them what they value most, but they might not tell me the truth. Alternatively, I can propose a division that's fair based on possible scenarios.} \\
\\ 
\texttt{Let me think about possible value functions for the other player. They could value bells, dolls, or threads differently. For example, if they value threads highly, they might want more threads. If they value the doll more, they might prefer that. Since I have threads as 4 each, which is higher than bells (0) and doll (2), maybe I should consider giving them threads if they don't value them as much. But without knowing their values, it's tricky.} \\
\\ 
\texttt{The key is to make a proposal that allows both of us to get items we value, so the total is maximized. Let's consider possible allocations.} \\
\\ 
\texttt{If I take the doll (2 points) and one thread (4), that's 6 points. The other player could take the remaining bell (0) and thread (4), totaling 4. But maybe there's a better split. Alternatively, if I take two threads (8 points) and the other takes the doll (2) and a bell (0), total would be 8 + 2 = 10. Wait, but the other player's value for the doll and bell depends on their own function. If they value the doll higher, maybe that's better.} \\
\\ 
\texttt{But since I don't know their values, I need to make a proposal that's likely to be complementary. Maybe propose taking the doll and one thread, leaving the bells and another thread for them. Alternatively, take two threads and a bell, but that might not be optimal.} \\
\\ 
\texttt{Alternatively, maybe I should ask the other player what they prefer. For example, I could message them:} \\
            }
        }
    }
    & & \\ \\

    \theutterance \stepcounter{utterance}  
    & & \multicolumn{4}{p{0.6\linewidth}}{
        \cellcolor[rgb]{0.9,0.9,0.9}{
            \makecell[{{p{\linewidth}}}]{
                \texttt{\tiny{[P1$\langle$GM]}}
                \texttt{<think>} \\
\texttt{Okay, let's see. I need to figure out how to divide the items between me and the other player to maximize the total score. The items are 2 bells, 1 doll, and 2 threads. My value function is bell: 0, doll: 2, thread: 4. The other player's values are secret, so I don't know what they value. The goal is to maximize the sum of both our scores.} \\
\\ 
\texttt{First, I should consider what the other player might value. Since I don't have their value function, I need to make some assumptions or ask them. But since I can only send 5 messages, I need to be efficient. Maybe I can ask them what they value most, but they might not tell me the truth. Alternatively, I can propose a division that's fair based on possible scenarios.} \\
\\ 
\texttt{Let me think about possible value functions for the other player. They could value bells, dolls, or threads differently. For example, if they value threads highly, they might want more threads. If they value the doll more, they might prefer that. Since I have threads as 4 each, which is higher than bells (0) and doll (2), maybe I should consider giving them threads if they don't value them as much. But without knowing their values, it's tricky.} \\
\\ 
\texttt{The key is to make a proposal that allows both of us to get items we value, so the total is maximized. Let's consider possible allocations.} \\
\\ 
\texttt{If I take the doll (2 points) and one thread (4), that's 6 points. The other player could take the remaining bell (0) and thread (4), totaling 4. But maybe there's a better split. Alternatively, if I take two threads (8 points) and the other takes the doll (2) and a bell (0), total would be 8 + 2 = 10. Wait, but the other player's value for the doll and bell depends on their own function. If they value the doll higher, maybe that's better.} \\
\\ 
\texttt{But since I don't know their values, I need to make a proposal that's likely to be complementary. Maybe propose taking the doll and one thread, leaving the bells and another thread for them. Alternatively, take two threads and a bell, but that might not be optimal.} \\
\\ 
\texttt{Alternatively, maybe I should ask the other player what they prefer. For example, I could message them:} \\
            }
        }
    }
    & \\ \\

    \theutterance \stepcounter{utterance}  
    & \multicolumn{4}{p{0.6\linewidth}}{
        \cellcolor[rgb]{0.9,0.9,0.9}{
            \makecell[{{p{\linewidth}}}]{
                \texttt{\tiny{[P1$\rangle$GM]}}
                \texttt{<think>} \\
\texttt{Okay, let's see. I need to figure out how to divide the items between me and the other player to maximize the total score. The items are 2 bells, 1 doll, and 2 threads. My value function is bell: 0, doll: 2, thread: 4. The other player's values are secret, so I don't know what they value. The goal is to maximize the sum of both our scores.} \\
\\ 
\texttt{First, I should consider what the other player might value. Since I don't have their value function, I need to make some assumptions or ask them. But since I can only send 5 messages, I need to be efficient. Maybe I can ask them what they value most, but they might not tell me the truth. Alternatively, I can propose a division that's fair based on possible scenarios.} \\
\\ 
\texttt{Let me think about possible value functions for the other player. They could value bells, dolls, or threads differently. For example, if they value threads highly, they might want more threads. If they value the doll more, they might prefer that. Since I have threads as 4 each, which is higher than bells (0) and doll (2), maybe I should consider giving them threads if they don't value them as much. But without knowing their values, it's tricky.} \\
\\ 
\texttt{The key is to make a proposal that allows both of us to get items we value, so the total is maximized. Let's consider possible allocations.} \\
\\ 
\texttt{If I take the doll (2 points) and one thread (4), that's 6 points. The other player could take the remaining bell (0) and thread (4), totaling 4. But maybe there's a better split. Alternatively, if I take two threads (8 points) and the other takes the doll (2) and a bell (0), total would be 8 + 2 = 10. Wait, but the other player's value for the doll and bell depends on their own function. If they value the doll higher, maybe that's better.} \\
\\ 
\texttt{But since I don't know their values, I need to make a proposal that's likely to be complementary. Maybe propose taking the doll and one thread, leaving the bells and another thread for them. Alternatively, take two threads and a bell, but that might not be optimal.} \\
\\ 
\texttt{Alternatively, maybe I should ask the other player what they prefer. For example, I could message them:} \\
            }
        }
    }
    & & \\ \\

    \theutterance \stepcounter{utterance}  
    & & & \multicolumn{2}{p{0.3\linewidth}}{
        \cellcolor[rgb]{0.9,0.9,0.9}{
            \makecell[{{p{\linewidth}}}]{
                \texttt{\tiny{[GM$|$GM]}}
                \texttt{continue} \\
            }
        }
    }
    & & \\ \\

    \theutterance \stepcounter{utterance}  
    & & & \multicolumn{2}{p{0.3\linewidth}}{
        \cellcolor[rgb]{0.9,0.9,0.9}{
            \makecell[{{p{\linewidth}}}]{
                \texttt{\tiny{[GM$|$GM]}}
                \texttt{<think>} \\
\texttt{Okay, let's see. I need to figure out how to divide the items between me and the other player to maximize the total score. The items are 2 bells, 1 doll, and 2 threads. My value function is bell: 0, doll: 2, thread: 4. The other player's values are secret, so I don't know what they value. The goal is to maximize the sum of both our scores.} \\
\\ 
\texttt{First, I should consider what the other player might value. Since I don't have their value function, I need to make some assumptions or ask them. But since I can only send 5 messages, I need to be efficient. Maybe I can ask them what they value most, but they might not tell me the truth. Alternatively, I can propose a division that's fair based on possible scenarios.} \\
\\ 
\texttt{Let me think about possible value functions for the other player. They could value bells, dolls, or threads differently. For example, if they value threads highly, they might want more threads. If they value the doll more, they might prefer that. Since I have threads as 4 each, which is higher than bells (0) and doll (2), maybe I should consider giving them threads if they don't value them as much. But without knowing their values, it's tricky.} \\
\\ 
\texttt{The key is to make a proposal that allows both of us to get items we value, so the total is maximized. Let's consider possible allocations.} \\
\\ 
\texttt{If I take the doll (2 points) and one thread (4), that's 6 points. The other player could take the remaining bell (0) and thread (4), totaling 4. But maybe there's a better split. Alternatively, if I take two threads (8 points) and the other takes the doll (2) and a bell (0), total would be 8 + 2 = 10. Wait, but the other player's value for the doll and bell depends on their own function. If they value the doll higher, maybe that's better.} \\
\\ 
\texttt{But since I don't know their values, I need to make a proposal that's likely to be complementary. Maybe propose taking the doll and one thread, leaving the bells and another thread for them. Alternatively, take two threads and a bell, but that might not be optimal.} \\
\\ 
\texttt{Alternatively, maybe I should ask the other player what they prefer. For example, I could message them:} \\
            }
        }
    }
    & & \\ \\

    \theutterance \stepcounter{utterance}  
    & & \multicolumn{4}{p{0.6\linewidth}}{
        \cellcolor[rgb]{0.9,0.9,0.9}{
            \makecell[{{p{\linewidth}}}]{
                \texttt{\tiny{[GM$\rangle$P2]}}
                \texttt{<think>} \\
\texttt{Okay, let's see. I need to figure out how to divide the items between me and the other player to maximize the total score. The items are 2 bells, 1 doll, and 2 threads. My value function is bell: 0, doll: 2, thread: 4. The other player's values are secret, so I don't know what they value. The goal is to maximize the sum of both our scores.} \\
\\ 
\texttt{First, I should consider what the other player might value. Since I don't have their value function, I need to make some assumptions or ask them. But since I can only send 5 messages, I need to be efficient. Maybe I can ask them what they value most, but they might not tell me the truth. Alternatively, I can propose a division that's fair based on possible scenarios.} \\
\\ 
\texttt{Let me think about possible value functions for the other player. They could value bells, dolls, or threads differently. For example, if they value threads highly, they might want more threads. If they value the doll more, they might prefer that. Since I have threads as 4 each, which is higher than bells (0) and doll (2), maybe I should consider giving them threads if they don't value them as much. But without knowing their values, it's tricky.} \\
\\ 
\texttt{The key is to make a proposal that allows both of us to get items we value, so the total is maximized. Let's consider possible allocations.} \\
\\ 
\texttt{If I take the doll (2 points) and one thread (4), that's 6 points. The other player could take the remaining bell (0) and thread (4), totaling 4. But maybe there's a better split. Alternatively, if I take two threads (8 points) and the other takes the doll (2) and a bell (0), total would be 8 + 2 = 10. Wait, but the other player's value for the doll and bell depends on their own function. If they value the doll higher, maybe that's better.} \\
\\ 
\texttt{But since I don't know their values, I need to make a proposal that's likely to be complementary. Maybe propose taking the doll and one thread, leaving the bells and another thread for them. Alternatively, take two threads and a bell, but that might not be optimal.} \\
\\ 
\texttt{Alternatively, maybe I should ask the other player what they prefer. For example, I could message them:} \\
            }
        }
    }
    & \\ \\

    \theutterance \stepcounter{utterance}  
    & & & \multicolumn{4}{p{0.6\linewidth}}{
        \cellcolor[rgb]{0.9,0.9,0.9}{
            \makecell[{{p{\linewidth}}}]{
                \texttt{\tiny{[GM$\langle$P2]}}
                \texttt{<think>} \\
\texttt{Okay, let's see. I need to figure out how to divide the items between me and the other player to maximize the total score. The items are 2 bells, 1 doll, and 2 threads. My value function is bell: 0, doll: 2, thread: 4. The other player's values are secret, so I don't know what they value. The goal is to maximize the sum of both our scores.} \\
\\ 
\texttt{First, I should consider what the other player might value. Since I don't have their value function, I need to make some assumptions or ask them. But since I can only send 5 messages, I need to be efficient. Maybe I can ask them what they value most, but they might not tell me the truth. Alternatively, I can propose a division that's fair based on possible scenarios.} \\
\\ 
\texttt{Let me think about possible value functions for the other player. They could value bells, dolls, or threads differently. For example, if they value threads highly, they might want more threads. If they value the doll more, they might prefer that. Since I have threads as 4 each, which is higher than bells (0) and doll (2), maybe I should consider giving them threads if they don't value them as much. But without knowing their values, it's tricky.} \\
\\ 
\texttt{The key is to make a proposal that allows both of us to get items we value, so the total is maximized. Let's consider possible allocations.} \\
\\ 
\texttt{If I take the doll (2 points) and one thread (4), that's 6 points. The other player could take the remaining bell (0) and thread (4), totaling 4. But maybe there's a better split. Alternatively, if I take two threads (8 points) and the other takes the doll (2) and a bell (0), total would be 8 + 2 = 10. Wait, but the other player's value for the doll and bell depends on their own function. If they value the doll higher, maybe that's better.} \\
\\ 
\texttt{But since I don't know their values, I need to make a proposal that's likely to be complementary. Maybe propose taking the doll and one thread, leaving the bells and another thread for them. Alternatively, take two threads and a bell, but that might not be optimal.} \\
\\ 
\texttt{Alternatively, maybe I should ask the other player what they prefer. For example, I could message them:} \\
            }
        }
    }
     \\ \\

    \theutterance \stepcounter{utterance}  
    & & & \multicolumn{2}{p{0.3\linewidth}}{
        \cellcolor[rgb]{0.9,0.9,0.9}{
            \makecell[{{p{\linewidth}}}]{
                \texttt{\tiny{[GM$|$GM]}}
                \texttt{continue} \\
            }
        }
    }
    & & \\ \\

    \theutterance \stepcounter{utterance}  
    & & & \multicolumn{2}{p{0.3\linewidth}}{
        \cellcolor[rgb]{0.9,0.9,0.9}{
            \makecell[{{p{\linewidth}}}]{
                \texttt{\tiny{[GM$|$GM]}}
                \texttt{<think>} \\
\texttt{Okay, let's see. I need to figure out how to divide the items between me and the other player to maximize the total score. The items are 2 bells, 1 doll, and 2 threads. My value function is bell: 0, doll: 2, thread: 4. The other player's values are secret, so I don't know what they value. The goal is to maximize the sum of both our scores.} \\
\\ 
\texttt{First, I should consider what the other player might value. Since I don't have their value function, I need to make some assumptions or ask them. But since I can only send 5 messages, I need to be efficient. Maybe I can ask them what they value most, but they might not tell me the truth. Alternatively, I can propose a division that's fair based on possible scenarios.} \\
\\ 
\texttt{Let me think about possible value functions for the other player. They could value bells, dolls, or threads differently. For example, if they value threads highly, they might want more threads. If they value the doll more, they might prefer that. Since I have threads as 4 each, which is higher than bells (0) and doll (2), maybe I should consider giving them threads if they don't value them as much. But without knowing their values, it's tricky.} \\
\\ 
\texttt{The key is to make a proposal that allows both of us to get items we value, so the total is maximized. Let's consider possible allocations.} \\
\\ 
\texttt{If I take the doll (2 points) and one thread (4), that's 6 points. The other player could take the remaining bell (0) and thread (4), totaling 4. But maybe there's a better split. Alternatively, if I take two threads (8 points) and the other takes the doll (2) and a bell (0), total would be 8 + 2 = 10. Wait, but the other player's value for the doll and bell depends on their own function. If they value the doll higher, maybe that's better.} \\
\\ 
\texttt{But since I don't know their values, I need to make a proposal that's likely to be complementary. Maybe propose taking the doll and one thread, leaving the bells and another thread for them. Alternatively, take two threads and a bell, but that might not be optimal.} \\
\\ 
\texttt{Alternatively, maybe I should ask the other player what they prefer. For example, I could message them:} \\
            }
        }
    }
    & & \\ \\

    \theutterance \stepcounter{utterance}  
    & & \multicolumn{4}{p{0.6\linewidth}}{
        \cellcolor[rgb]{0.9,0.9,0.9}{
            \makecell[{{p{\linewidth}}}]{
                \texttt{\tiny{[P1$\langle$GM]}}
                \texttt{<think>} \\
\texttt{Okay, let's see. I need to figure out how to divide the items between me and the other player to maximize the total score. The items are 2 bells, 1 doll, and 2 threads. My value function is bell: 0, doll: 2, thread: 4. The other player's values are secret, so I don't know what they value. The goal is to maximize the sum of both our scores.} \\
\\ 
\texttt{First, I should consider what the other player might value. Since I don't have their value function, I need to make some assumptions or ask them. But since I can only send 5 messages, I need to be efficient. Maybe I can ask them what they value most, but they might not tell me the truth. Alternatively, I can propose a division that's fair based on possible scenarios.} \\
\\ 
\texttt{Let me think about possible value functions for the other player. They could value bells, dolls, or threads differently. For example, if they value threads highly, they might want more threads. If they value the doll more, they might prefer that. Since I have threads as 4 each, which is higher than bells (0) and doll (2), maybe I should consider giving them threads if they don't value them as much. But without knowing their values, it's tricky.} \\
\\ 
\texttt{The key is to make a proposal that allows both of us to get items we value, so the total is maximized. Let's consider possible allocations.} \\
\\ 
\texttt{If I take the doll (2 points) and one thread (4), that's 6 points. The other player could take the remaining bell (0) and thread (4), totaling 4. But maybe there's a better split. Alternatively, if I take two threads (8 points) and the other takes the doll (2) and a bell (0), total would be 8 + 2 = 10. Wait, but the other player's value for the doll and bell depends on their own function. If they value the doll higher, maybe that's better.} \\
\\ 
\texttt{But since I don't know their values, I need to make a proposal that's likely to be complementary. Maybe propose taking the doll and one thread, leaving the bells and another thread for them. Alternatively, take two threads and a bell, but that might not be optimal.} \\
\\ 
\texttt{Alternatively, maybe I should ask the other player what they prefer. For example, I could message them:} \\
            }
        }
    }
    & \\ \\

    \theutterance \stepcounter{utterance}  
    & \multicolumn{4}{p{0.6\linewidth}}{
        \cellcolor[rgb]{0.9,0.9,0.9}{
            \makecell[{{p{\linewidth}}}]{
                \texttt{\tiny{[P1$\rangle$GM]}}
                \texttt{<think>} \\
\texttt{Okay, let's see. I need to figure out how to divide the items between me and the other player to maximize the total score. The items are 2 bells, 1 doll, and 2 threads. My value function is bell: 0, doll: 2, thread: 4. The other player's values are secret, so I don't know what they value. The goal is to maximize the sum of both our scores.} \\
\\ 
\texttt{First, I should consider what the other player might value. Since I don't have their value function, I need to make some assumptions or ask them. But since I can only send 5 messages, I need to be efficient. Maybe I can ask them what they value most, but they might not tell me the truth. Alternatively, I can propose a division that's fair based on possible scenarios.} \\
\\ 
\texttt{Let me think about possible value functions for the other player. They could value bells, dolls, or threads differently. For example, if they value threads highly, they might want more threads. If they value the doll more, they might prefer that. Since I have threads as 4 each, which is higher than bells (0) and doll (2), maybe I should consider giving them threads if they don't value them as much. But without knowing their values, it's tricky.} \\
\\ 
\texttt{The key is to make a proposal that allows both of us to get items we value, so the total is maximized. Let's consider possible allocations.} \\
\\ 
\texttt{If I take the doll (2 points) and one thread (4), that's 6 points. The other player could take the remaining bell (0) and thread (4), totaling 4. But maybe there's a better split. Alternatively, if I take two threads (8 points) and the other takes the doll (2) and a bell (0), total would be 8 + 2 = 10. Wait, but the other player's value for the doll and bell depends on their own function. If they value the doll higher, maybe that's better.} \\
\\ 
\texttt{But since I don't know their values, I need to make a proposal that's likely to be complementary. Maybe propose taking the doll and one thread, leaving the bells and another thread for them. Alternatively, take two threads and a bell, but that might not be optimal.} \\
\\ 
\texttt{Alternatively, maybe I should ask the other player what they prefer. For example, I could message them:} \\
            }
        }
    }
    & & \\ \\

    \theutterance \stepcounter{utterance}  
    & & & \multicolumn{2}{p{0.3\linewidth}}{
        \cellcolor[rgb]{0.9,0.9,0.9}{
            \makecell[{{p{\linewidth}}}]{
                \texttt{\tiny{[GM$|$GM]}}
                \texttt{continue} \\
            }
        }
    }
    & & \\ \\

    \theutterance \stepcounter{utterance}  
    & & & \multicolumn{2}{p{0.3\linewidth}}{
        \cellcolor[rgb]{0.9,0.9,0.9}{
            \makecell[{{p{\linewidth}}}]{
                \texttt{\tiny{[GM$|$GM]}}
                \texttt{<think>} \\
\texttt{Okay, let's see. I need to figure out how to divide the items between me and the other player to maximize the total score. The items are 2 bells, 1 doll, and 2 threads. My value function is bell: 0, doll: 2, thread: 4. The other player's values are secret, so I don't know what they value. The goal is to maximize the sum of both our scores.} \\
\\ 
\texttt{First, I should consider what the other player might value. Since I don't have their value function, I need to make some assumptions or ask them. But since I can only send 5 messages, I need to be efficient. Maybe I can ask them what they value most, but they might not tell me the truth. Alternatively, I can propose a division that's fair based on possible scenarios.} \\
\\ 
\texttt{Let me think about possible value functions for the other player. They could value bells, dolls, or threads differently. For example, if they value threads highly, they might want more threads. If they value the doll more, they might prefer that. Since I have threads as 4 each, which is higher than bells (0) and doll (2), maybe I should consider giving them threads if they don't value them as much. But without knowing their values, it's tricky.} \\
\\ 
\texttt{The key is to make a proposal that allows both of us to get items we value, so the total is maximized. Let's consider possible allocations.} \\
\\ 
\texttt{If I take the doll (2 points) and one thread (4), that's 6 points. The other player could take the remaining bell (0) and thread (4), totaling 4. But maybe there's a better split. Alternatively, if I take two threads (8 points) and the other takes the doll (2) and a bell (0), total would be 8 + 2 = 10. Wait, but the other player's value for the doll and bell depends on their own function. If they value the doll higher, maybe that's better.} \\
\\ 
\texttt{But since I don't know their values, I need to make a proposal that's likely to be complementary. Maybe propose taking the doll and one thread, leaving the bells and another thread for them. Alternatively, take two threads and a bell, but that might not be optimal.} \\
\\ 
\texttt{Alternatively, maybe I should ask the other player what they prefer. For example, I could message them:} \\
            }
        }
    }
    & & \\ \\

    \theutterance \stepcounter{utterance}  
    & & \multicolumn{4}{p{0.6\linewidth}}{
        \cellcolor[rgb]{0.9,0.9,0.9}{
            \makecell[{{p{\linewidth}}}]{
                \texttt{\tiny{[GM$\rangle$P2]}}
                \texttt{<think>} \\
\texttt{Okay, let's see. I need to figure out how to divide the items between me and the other player to maximize the total score. The items are 2 bells, 1 doll, and 2 threads. My value function is bell: 0, doll: 2, thread: 4. The other player's values are secret, so I don't know what they value. The goal is to maximize the sum of both our scores.} \\
\\ 
\texttt{First, I should consider what the other player might value. Since I don't have their value function, I need to make some assumptions or ask them. But since I can only send 5 messages, I need to be efficient. Maybe I can ask them what they value most, but they might not tell me the truth. Alternatively, I can propose a division that's fair based on possible scenarios.} \\
\\ 
\texttt{Let me think about possible value functions for the other player. They could value bells, dolls, or threads differently. For example, if they value threads highly, they might want more threads. If they value the doll more, they might prefer that. Since I have threads as 4 each, which is higher than bells (0) and doll (2), maybe I should consider giving them threads if they don't value them as much. But without knowing their values, it's tricky.} \\
\\ 
\texttt{The key is to make a proposal that allows both of us to get items we value, so the total is maximized. Let's consider possible allocations.} \\
\\ 
\texttt{If I take the doll (2 points) and one thread (4), that's 6 points. The other player could take the remaining bell (0) and thread (4), totaling 4. But maybe there's a better split. Alternatively, if I take two threads (8 points) and the other takes the doll (2) and a bell (0), total would be 8 + 2 = 10. Wait, but the other player's value for the doll and bell depends on their own function. If they value the doll higher, maybe that's better.} \\
\\ 
\texttt{But since I don't know their values, I need to make a proposal that's likely to be complementary. Maybe propose taking the doll and one thread, leaving the bells and another thread for them. Alternatively, take two threads and a bell, but that might not be optimal.} \\
\\ 
\texttt{Alternatively, maybe I should ask the other player what they prefer. For example, I could message them:} \\
            }
        }
    }
    & \\ \\

    \theutterance \stepcounter{utterance}  
    & & & \multicolumn{4}{p{0.6\linewidth}}{
        \cellcolor[rgb]{0.9,0.9,0.9}{
            \makecell[{{p{\linewidth}}}]{
                \texttt{\tiny{[GM$\langle$P2]}}
                \texttt{<think>} \\
\texttt{Okay, let's see. I need to figure out how to divide the items between me and the other player to maximize the total score. The items are 2 bells, 1 doll, and 2 threads. My value function is bell: 0, doll: 2, thread: 4. The other player's values are secret, so I don't know what they value. The goal is to maximize the sum of both our scores.} \\
\\ 
\texttt{First, I should consider what the other player might value. Since I don't have their value function, I need to make some assumptions or ask them. But since I can only send 5 messages, I need to be efficient. Maybe I can ask them what they value most, but they might not tell me the truth. Alternatively, I can propose a division that's fair based on possible scenarios.} \\
\\ 
\texttt{Let me think about possible value functions for the other player. They could value bells, dolls, or threads differently. For example, if they value threads highly, they might want more threads. If they value the doll more, they might prefer that. Since I have threads as 4 each, which is higher than bells (0) and doll (2), maybe I should consider giving them threads if they don't value them as much. But without knowing their values, it's tricky.} \\
\\ 
\texttt{The key is to make a proposal that allows both of us to get items we value, so the total is maximized. Let's consider possible allocations.} \\
\\ 
\texttt{If I take the doll (2 points) and one thread (4), that's 6 points. The other player could take the remaining bell (0) and thread (4), totaling 4. But maybe there's a better split. Alternatively, if I take two threads (8 points) and the other takes the doll (2) and a bell (0), total would be 8 + 2 = 10. Wait, but the other player's value for the doll and bell depends on their own function. If they value the doll higher, maybe that's better.} \\
\\ 
\texttt{But since I don't know their values, I need to make a proposal that's likely to be complementary. Maybe propose taking the doll and one thread, leaving the bells and another thread for them. Alternatively, take two threads and a bell, but that might not be optimal.} \\
\\ 
\texttt{Alternatively, maybe I should ask the other player what they prefer. For example, I could message them:} \\
            }
        }
    }
     \\ \\

    \theutterance \stepcounter{utterance}  
    & & & \multicolumn{2}{p{0.3\linewidth}}{
        \cellcolor[rgb]{0.9,0.9,0.9}{
            \makecell[{{p{\linewidth}}}]{
                \texttt{\tiny{[GM$|$GM]}}
                \texttt{continue} \\
            }
        }
    }
    & & \\ \\

    \theutterance \stepcounter{utterance}  
    & & & \multicolumn{2}{p{0.3\linewidth}}{
        \cellcolor[rgb]{0.9,0.9,0.9}{
            \makecell[{{p{\linewidth}}}]{
                \texttt{\tiny{[GM$|$GM]}}
                \texttt{<think>} \\
\texttt{Okay, let's see. I need to figure out how to divide the items between me and the other player to maximize the total score. The items are 2 bells, 1 doll, and 2 threads. My value function is bell: 0, doll: 2, thread: 4. The other player's values are secret, so I don't know what they value. The goal is to maximize the sum of both our scores.} \\
\\ 
\texttt{First, I should consider what the other player might value. Since I don't have their value function, I need to make some assumptions or ask them. But since I can only send 5 messages, I need to be efficient. Maybe I can ask them what they value most, but they might not tell me the truth. Alternatively, I can propose a division that's fair based on possible scenarios.} \\
\\ 
\texttt{Let me think about possible value functions for the other player. They could value bells, dolls, or threads differently. For example, if they value threads highly, they might want more threads. If they value the doll more, they might prefer that. Since I have threads as 4 each, which is higher than bells (0) and doll (2), maybe I should consider giving them threads if they don't value them as much. But without knowing their values, it's tricky.} \\
\\ 
\texttt{The key is to make a proposal that allows both of us to get items we value, so the total is maximized. Let's consider possible allocations.} \\
\\ 
\texttt{If I take the doll (2 points) and one thread (4), that's 6 points. The other player could take the remaining bell (0) and thread (4), totaling 4. But maybe there's a better split. Alternatively, if I take two threads (8 points) and the other takes the doll (2) and a bell (0), total would be 8 + 2 = 10. Wait, but the other player's value for the doll and bell depends on their own function. If they value the doll higher, maybe that's better.} \\
\\ 
\texttt{But since I don't know their values, I need to make a proposal that's likely to be complementary. Maybe propose taking the doll and one thread, leaving the bells and another thread for them. Alternatively, take two threads and a bell, but that might not be optimal.} \\
\\ 
\texttt{Alternatively, maybe I should ask the other player what they prefer. For example, I could message them:} \\
            }
        }
    }
    & & \\ \\

    \theutterance \stepcounter{utterance}  
    & & \multicolumn{4}{p{0.6\linewidth}}{
        \cellcolor[rgb]{0.9,0.9,0.9}{
            \makecell[{{p{\linewidth}}}]{
                \texttt{\tiny{[P1$\langle$GM]}}
                \texttt{<think>} \\
\texttt{Okay, let's see. I need to figure out how to divide the items between me and the other player to maximize the total score. The items are 2 bells, 1 doll, and 2 threads. My value function is bell: 0, doll: 2, thread: 4. The other player's values are secret, so I don't know what they value. The goal is to maximize the sum of both our scores.} \\
\\ 
\texttt{First, I should consider what the other player might value. Since I don't have their value function, I need to make some assumptions or ask them. But since I can only send 5 messages, I need to be efficient. Maybe I can ask them what they value most, but they might not tell me the truth. Alternatively, I can propose a division that's fair based on possible scenarios.} \\
\\ 
\texttt{Let me think about possible value functions for the other player. They could value bells, dolls, or threads differently. For example, if they value threads highly, they might want more threads. If they value the doll more, they might prefer that. Since I have threads as 4 each, which is higher than bells (0) and doll (2), maybe I should consider giving them threads if they don't value them as much. But without knowing their values, it's tricky.} \\
\\ 
\texttt{The key is to make a proposal that allows both of us to get items we value, so the total is maximized. Let's consider possible allocations.} \\
\\ 
\texttt{If I take the doll (2 points) and one thread (4), that's 6 points. The other player could take the remaining bell (0) and thread (4), totaling 4. But maybe there's a better split. Alternatively, if I take two threads (8 points) and the other takes the doll (2) and a bell (0), total would be 8 + 2 = 10. Wait, but the other player's value for the doll and bell depends on their own function. If they value the doll higher, maybe that's better.} \\
\\ 
\texttt{But since I don't know their values, I need to make a proposal that's likely to be complementary. Maybe propose taking the doll and one thread, leaving the bells and another thread for them. Alternatively, take two threads and a bell, but that might not be optimal.} \\
\\ 
\texttt{Alternatively, maybe I should ask the other player what they prefer. For example, I could message them:} \\
            }
        }
    }
    & \\ \\

    \theutterance \stepcounter{utterance}  
    & \multicolumn{4}{p{0.6\linewidth}}{
        \cellcolor[rgb]{0.9,0.9,0.9}{
            \makecell[{{p{\linewidth}}}]{
                \texttt{\tiny{[P1$\rangle$GM]}}
                \texttt{<think>} \\
\texttt{Okay, let's see. I need to figure out how to divide the items between me and the other player to maximize the total score. The items are 2 bells, 1 doll, and 2 threads. My value function is bell: 0, doll: 2, thread: 4. The other player's values are secret, so I don't know what they value. The goal is to maximize the sum of both our scores.} \\
\\ 
\texttt{First, I should consider what the other player might value. Since I don't have their value function, I need to make some assumptions or ask them. But since I can only send 5 messages, I need to be efficient. Maybe I can ask them what they value most, but they might not tell me the truth. Alternatively, I can propose a division that's fair based on possible scenarios.} \\
\\ 
\texttt{Let me think about possible value functions for the other player. They could value bells, dolls, or threads differently. For example, if they value threads highly, they might want more threads. If they value the doll more, they might prefer that. Since I have threads as 4 each, which is higher than bells (0) and doll (2), maybe I should consider giving them threads if they don't value them as much. But without knowing their values, it's tricky.} \\
\\ 
\texttt{The key is to make a proposal that allows both of us to get items we value, so the total is maximized. Let's consider possible allocations.} \\
\\ 
\texttt{If I take the doll (2 points) and one thread (4), that's 6 points. The other player could take the remaining bell (0) and thread (4), totaling 4. But maybe there's a better split. Alternatively, if I take two threads (8 points) and the other takes the doll (2) and a bell (0), total would be 8 + 2 = 10. Wait, but the other player's value for the doll and bell depends on their own function. If they value the doll higher, maybe that's better.} \\
\\ 
\texttt{But since I don't know their values, I need to make a proposal that's likely to be complementary. Maybe propose taking the doll and one thread, leaving the bells and another thread for them. Alternatively, take two threads and a bell, but that might not be optimal.} \\
\\ 
\texttt{Alternatively, maybe I should ask the other player what they prefer. For example, I could message them:} \\
            }
        }
    }
    & & \\ \\

    \theutterance \stepcounter{utterance}  
    & & & \multicolumn{2}{p{0.3\linewidth}}{
        \cellcolor[rgb]{0.9,0.9,0.9}{
            \makecell[{{p{\linewidth}}}]{
                \texttt{\tiny{[GM$|$GM]}}
                \texttt{continue} \\
            }
        }
    }
    & & \\ \\

    \theutterance \stepcounter{utterance}  
    & & & \multicolumn{2}{p{0.3\linewidth}}{
        \cellcolor[rgb]{0.9,0.9,0.9}{
            \makecell[{{p{\linewidth}}}]{
                \texttt{\tiny{[GM$|$GM]}}
                \texttt{<think>} \\
\texttt{Okay, let's see. I need to figure out how to divide the items between me and the other player to maximize the total score. The items are 2 bells, 1 doll, and 2 threads. My value function is bell: 0, doll: 2, thread: 4. The other player's values are secret, so I don't know what they value. The goal is to maximize the sum of both our scores.} \\
\\ 
\texttt{First, I should consider what the other player might value. Since I don't have their value function, I need to make some assumptions or ask them. But since I can only send 5 messages, I need to be efficient. Maybe I can ask them what they value most, but they might not tell me the truth. Alternatively, I can propose a division that's fair based on possible scenarios.} \\
\\ 
\texttt{Let me think about possible value functions for the other player. They could value bells, dolls, or threads differently. For example, if they value threads highly, they might want more threads. If they value the doll more, they might prefer that. Since I have threads as 4 each, which is higher than bells (0) and doll (2), maybe I should consider giving them threads if they don't value them as much. But without knowing their values, it's tricky.} \\
\\ 
\texttt{The key is to make a proposal that allows both of us to get items we value, so the total is maximized. Let's consider possible allocations.} \\
\\ 
\texttt{If I take the doll (2 points) and one thread (4), that's 6 points. The other player could take the remaining bell (0) and thread (4), totaling 4. But maybe there's a better split. Alternatively, if I take two threads (8 points) and the other takes the doll (2) and a bell (0), total would be 8 + 2 = 10. Wait, but the other player's value for the doll and bell depends on their own function. If they value the doll higher, maybe that's better.} \\
\\ 
\texttt{But since I don't know their values, I need to make a proposal that's likely to be complementary. Maybe propose taking the doll and one thread, leaving the bells and another thread for them. Alternatively, take two threads and a bell, but that might not be optimal.} \\
\\ 
\texttt{Alternatively, maybe I should ask the other player what they prefer. For example, I could message them:} \\
            }
        }
    }
    & & \\ \\

    \theutterance \stepcounter{utterance}  
    & & \multicolumn{4}{p{0.6\linewidth}}{
        \cellcolor[rgb]{0.9,0.9,0.9}{
            \makecell[{{p{\linewidth}}}]{
                \texttt{\tiny{[GM$\rangle$P2]}}
                \texttt{<think>} \\
\texttt{Okay, let's see. I need to figure out how to divide the items between me and the other player to maximize the total score. The items are 2 bells, 1 doll, and 2 threads. My value function is bell: 0, doll: 2, thread: 4. The other player's values are secret, so I don't know what they value. The goal is to maximize the sum of both our scores.} \\
\\ 
\texttt{First, I should consider what the other player might value. Since I don't have their value function, I need to make some assumptions or ask them. But since I can only send 5 messages, I need to be efficient. Maybe I can ask them what they value most, but they might not tell me the truth. Alternatively, I can propose a division that's fair based on possible scenarios.} \\
\\ 
\texttt{Let me think about possible value functions for the other player. They could value bells, dolls, or threads differently. For example, if they value threads highly, they might want more threads. If they value the doll more, they might prefer that. Since I have threads as 4 each, which is higher than bells (0) and doll (2), maybe I should consider giving them threads if they don't value them as much. But without knowing their values, it's tricky.} \\
\\ 
\texttt{The key is to make a proposal that allows both of us to get items we value, so the total is maximized. Let's consider possible allocations.} \\
\\ 
\texttt{If I take the doll (2 points) and one thread (4), that's 6 points. The other player could take the remaining bell (0) and thread (4), totaling 4. But maybe there's a better split. Alternatively, if I take two threads (8 points) and the other takes the doll (2) and a bell (0), total would be 8 + 2 = 10. Wait, but the other player's value for the doll and bell depends on their own function. If they value the doll higher, maybe that's better.} \\
\\ 
\texttt{But since I don't know their values, I need to make a proposal that's likely to be complementary. Maybe propose taking the doll and one thread, leaving the bells and another thread for them. Alternatively, take two threads and a bell, but that might not be optimal.} \\
\\ 
\texttt{Alternatively, maybe I should ask the other player what they prefer. For example, I could message them:} \\
            }
        }
    }
    & \\ \\

    \theutterance \stepcounter{utterance}  
    & & & \multicolumn{4}{p{0.6\linewidth}}{
        \cellcolor[rgb]{0.9,0.9,0.9}{
            \makecell[{{p{\linewidth}}}]{
                \texttt{\tiny{[GM$\langle$P2]}}
                \texttt{<think>} \\
\texttt{Okay, let's see. I need to figure out how to divide the items between me and the other player to maximize the total score. The items are 2 bells, 1 doll, and 2 threads. My value function is bell: 0, doll: 2, thread: 4. The other player's values are secret, so I don't know what they value. The goal is to maximize the sum of both our scores.} \\
\\ 
\texttt{First, I should consider what the other player might value. Since I don't have their value function, I need to make some assumptions or ask them. But since I can only send 5 messages, I need to be efficient. Maybe I can ask them what they value most, but they might not tell me the truth. Alternatively, I can propose a division that's fair based on possible scenarios.} \\
\\ 
\texttt{Let me think about possible value functions for the other player. They could value bells, dolls, or threads differently. For example, if they value threads highly, they might want more threads. If they value the doll more, they might prefer that. Since I have threads as 4 each, which is higher than bells (0) and doll (2), maybe I should consider giving them threads if they don't value them as much. But without knowing their values, it's tricky.} \\
\\ 
\texttt{The key is to make a proposal that allows both of us to get items we value, so the total is maximized. Let's consider possible allocations.} \\
\\ 
\texttt{If I take the doll (2 points) and one thread (4), that's 6 points. The other player could take the remaining bell (0) and thread (4), totaling 4. But maybe there's a better split. Alternatively, if I take two threads (8 points) and the other takes the doll (2) and a bell (0), total would be 8 + 2 = 10. Wait, but the other player's value for the doll and bell depends on their own function. If they value the doll higher, maybe that's better.} \\
\\ 
\texttt{But since I don't know their values, I need to make a proposal that's likely to be complementary. Maybe propose taking the doll and one thread, leaving the bells and another thread for them. Alternatively, take two threads and a bell, but that might not be optimal.} \\
\\ 
\texttt{Alternatively, maybe I should ask the other player what they prefer. For example, I could message them:} \\
            }
        }
    }
     \\ \\

    \theutterance \stepcounter{utterance}  
    & & & \multicolumn{2}{p{0.3\linewidth}}{
        \cellcolor[rgb]{0.9,0.9,0.9}{
            \makecell[{{p{\linewidth}}}]{
                \texttt{\tiny{[GM$|$GM]}}
                \texttt{continue} \\
            }
        }
    }
    & & \\ \\

    \theutterance \stepcounter{utterance}  
    & & & \multicolumn{2}{p{0.3\linewidth}}{
        \cellcolor[rgb]{0.9,0.9,0.9}{
            \makecell[{{p{\linewidth}}}]{
                \texttt{\tiny{[GM$|$GM]}}
                \texttt{<think>} \\
\texttt{Okay, let's see. I need to figure out how to divide the items between me and the other player to maximize the total score. The items are 2 bells, 1 doll, and 2 threads. My value function is bell: 0, doll: 2, thread: 4. The other player's values are secret, so I don't know what they value. The goal is to maximize the sum of both our scores.} \\
\\ 
\texttt{First, I should consider what the other player might value. Since I don't have their value function, I need to make some assumptions or ask them. But since I can only send 5 messages, I need to be efficient. Maybe I can ask them what they value most, but they might not tell me the truth. Alternatively, I can propose a division that's fair based on possible scenarios.} \\
\\ 
\texttt{Let me think about possible value functions for the other player. They could value bells, dolls, or threads differently. For example, if they value threads highly, they might want more threads. If they value the doll more, they might prefer that. Since I have threads as 4 each, which is higher than bells (0) and doll (2), maybe I should consider giving them threads if they don't value them as much. But without knowing their values, it's tricky.} \\
\\ 
\texttt{The key is to make a proposal that allows both of us to get items we value, so the total is maximized. Let's consider possible allocations.} \\
\\ 
\texttt{If I take the doll (2 points) and one thread (4), that's 6 points. The other player could take the remaining bell (0) and thread (4), totaling 4. But maybe there's a better split. Alternatively, if I take two threads (8 points) and the other takes the doll (2) and a bell (0), total would be 8 + 2 = 10. Wait, but the other player's value for the doll and bell depends on their own function. If they value the doll higher, maybe that's better.} \\
\\ 
\texttt{But since I don't know their values, I need to make a proposal that's likely to be complementary. Maybe propose taking the doll and one thread, leaving the bells and another thread for them. Alternatively, take two threads and a bell, but that might not be optimal.} \\
\\ 
\texttt{Alternatively, maybe I should ask the other player what they prefer. For example, I could message them:} \\
            }
        }
    }
    & & \\ \\

    \theutterance \stepcounter{utterance}  
    & & \multicolumn{4}{p{0.6\linewidth}}{
        \cellcolor[rgb]{0.9,0.9,0.9}{
            \makecell[{{p{\linewidth}}}]{
                \texttt{\tiny{[P1$\langle$GM]}}
                \texttt{<think>} \\
\texttt{Okay, let's see. I need to figure out how to divide the items between me and the other player to maximize the total score. The items are 2 bells, 1 doll, and 2 threads. My value function is bell: 0, doll: 2, thread: 4. The other player's values are secret, so I don't know what they value. The goal is to maximize the sum of both our scores.} \\
\\ 
\texttt{First, I should consider what the other player might value. Since I don't have their value function, I need to make some assumptions or ask them. But since I can only send 5 messages, I need to be efficient. Maybe I can ask them what they value most, but they might not tell me the truth. Alternatively, I can propose a division that's fair based on possible scenarios.} \\
\\ 
\texttt{Let me think about possible value functions for the other player. They could value bells, dolls, or threads differently. For example, if they value threads highly, they might want more threads. If they value the doll more, they might prefer that. Since I have threads as 4 each, which is higher than bells (0) and doll (2), maybe I should consider giving them threads if they don't value them as much. But without knowing their values, it's tricky.} \\
\\ 
\texttt{The key is to make a proposal that allows both of us to get items we value, so the total is maximized. Let's consider possible allocations.} \\
\\ 
\texttt{If I take the doll (2 points) and one thread (4), that's 6 points. The other player could take the remaining bell (0) and thread (4), totaling 4. But maybe there's a better split. Alternatively, if I take two threads (8 points) and the other takes the doll (2) and a bell (0), total would be 8 + 2 = 10. Wait, but the other player's value for the doll and bell depends on their own function. If they value the doll higher, maybe that's better.} \\
\\ 
\texttt{But since I don't know their values, I need to make a proposal that's likely to be complementary. Maybe propose taking the doll and one thread, leaving the bells and another thread for them. Alternatively, take two threads and a bell, but that might not be optimal.} \\
\\ 
\texttt{Alternatively, maybe I should ask the other player what they prefer. For example, I could message them:} \\
\\ 
\\ 
\texttt{The time is up. It is now your turn to submit a secret proposal.} \\
            }
        }
    }
    & \\ \\

    \theutterance \stepcounter{utterance}  
    & \multicolumn{4}{p{0.6\linewidth}}{
        \cellcolor[rgb]{0.9,0.9,0.9}{
            \makecell[{{p{\linewidth}}}]{
                \texttt{\tiny{[P1$\rangle$GM]}}
                \texttt{<think>} \\
\texttt{Okay, let's see. I need to make a secret proposal here. The items are 2 bells, 1 doll, and 2 threads. My value function is bell: 0, doll: 2, thread: 4. The other player's values are unknown. The goal is to maximize the total score for both of us.} \\
\\ 
\texttt{Since I don't know their values, I should aim for a division that's likely to be complementary. If I take the items I value most, that would be the threads and the doll. But if I take both threads, that's 8 points for me, but the other player would get the doll and a bell. If their value for the doll is high, maybe that's okay. But if they value threads more, then taking two threads might not be the best.} \\
\\ 
\texttt{Alternatively, if I take one thread and the doll (totaling 6 points), the other player could get the remaining thread and both bells. If they value threads, they'd get 4 points from the thread. If they value bells, but my value for bells is 0, so maybe they don't value them. But the doll is 2 for me, so if they take the doll, maybe they value it higher. Hmm.} \\
\\ 
\texttt{Wait, the total maximum possible score if we both take what we value most would be if I take two threads (8) and they take the doll (assuming they value it). But if their value for the doll is higher than mine, that's good. But if they don't value the doll, then maybe they'd prefer threads. But since I don't know, maybe the safest bet is to take the two threads and the doll, but that might not leave anything for them. Wait, there are only 2 bells, 1 doll, 2 threads. If I take the doll and two threads, that's all the threads and the doll. The other player would get the two bells, which I value at 0, but maybe they value them. If they do, then their score would be based on bells. But if they don't, then maybe that's not good.} \\
\\ 
\texttt{Alternatively, if I take one thread and the doll, leaving one thread and two bells for them. That way, if they value threads, they get 4, and if they value bells, they get whatever their value is. But since I don't know, maybe this is a better split} \\
            }
        }
    }
    & & \\ \\

    \theutterance \stepcounter{utterance}  
    & & & \multicolumn{2}{p{0.3\linewidth}}{
        \cellcolor[rgb]{0.9,0.9,0.9}{
            \makecell[{{p{\linewidth}}}]{
                \texttt{\tiny{[GM$|$GM]}}
                \texttt{continue} \\
            }
        }
    }
    & & \\ \\

    \theutterance \stepcounter{utterance}  
    & & & \multicolumn{2}{p{0.3\linewidth}}{
        \cellcolor[rgb]{0.9,0.9,0.9}{
            \makecell[{{p{\linewidth}}}]{
                \texttt{\tiny{[GM$|$GM]}}
                \texttt{player was instructed to make a proposal} \\
            }
        }
    }
    & & \\ \\

\end{supertabular}
}

\end{document}
